\section{Results: natural distribution, fine-tuning and label robustness}\label{sec:robust}

\subsection{Unaltered test distributions}
Without any resampling, in-ROI artifacts are naturally associated with the label in every cohort (thyroid calipers:
37.6\% of benign vs 24.1\% of malignant nodules; in-lesion hair: 31\% of melanomas vs 20\% of benign lesions;
capsule debris on the lesion: 23\% of polyp-like lesions vs 10\% of erosions). We trained on natural training data
and tested on untouched data: the official TN3K test split, each ISIC source held out in turn (a real change of
hospital and device), and held-out capsule frame blocks. On the shortcut-conflicting pairs, masking \emph{lowers}
AUROC for thyroid ($-0.102$ [$-0.139, -0.066$]) and for two of the three held-out hospitals (BCN $-0.016$
[$-0.027, -0.005$], MSK $-0.056$ [$-0.084, -0.029$]); it helps for HAM ($+0.053$) and does not differ significantly
from ERM on capsule. U-MtE with group balancing improves on masking for these pairs in every cohort in which it was
evaluated ($+0.014$ to $+0.065$) and gives the best or tied-best cross-group AUROC in three of five settings (HAM
0.786, BCN 0.734 vs 0.735, capsule 0.947); plain group balancing is best on thyroid and MSK. Without artifact labels,
U-MtE improves on masking for thyroid ($+0.051$) and HAM ($+0.023$) only.

\subsection{End-to-end fine-tuning}\label{sec:ft}
Fine-tuning ResNet-50 end-to-end (five folds) reproduces the frozen-backbone findings. Capsule: masking harms
in-ROI ($-0.284$ [$-0.347, -0.226$]) and helps out-of-ROI ($+0.301$), crossover $+0.584$ [$+0.469, +0.691$].
Thyroid: the crossover points the same way ($+0.057$ [$-0.021, +0.130$], five folds), and an end-to-end analogue of
U-MtE---an invariance penalty between masked images with and without a generic overlay---improves on masking by
$+0.143$ [$+0.099, +0.188$] without shortcut flipping (reversed 0.679, correlated 0.703), matching the best
label-using arm. It fails on capsule ($-0.021$), where the artifact resembles the pathology. Ovary (1,202 images)
is underpowered: U-MtE end-to-end has the best reversed, correlated and clean AUROC, with CIs overlapping masking.
A fine-tuned ViT-S reproduces the crossover on thyroid ($+0.168$ [$+0.074, +0.265$]), but the end-to-end U-MtE
analogue does not improve on masking with this architecture ($+0.009$ [$-0.021, +0.040$]).

\subsection{Robustness to artifact-label definitions}
Artifact masks for calipers come from a detector (visual audit: 95\% of thyroid detections lie on real calipers;
92--97\% of caliper-free images contain no visible caliper). Re-running the thyroid and ovary analyses with
large-marker-only ($\ge 50$ px) and strict-location ($r\ge0.7$ / $r<0.05$) definitions leaves the crossover
(thyroid $+0.215$ to $+0.262$; ovary $+0.170$ to $+0.195$) and the U-MtE$_{\text{protect}}$ gain over masking
(thyroid $+0.217$ to $+0.282$; ovary $+0.043$ to $+0.067$) significant in every variant; both grow under stricter
definitions, as expected if detector errors dilute rather than create the effects.
